\documentclass[11pt,a4paper]{article}
\usepackage[utf8]{inputenc}
\usepackage[T1]{fontenc}
\usepackage{amsmath,amssymb,amsthm}
\usepackage{graphicx}
\usepackage{hyperref}
\usepackage{booktabs}
\usepackage{geometry}
\geometry{margin=1in}

\newtheorem{theorem}{Theorem}
\newtheorem{definition}[theorem]{Definition}

\hypersetup{
    pdftitle={Foundation: One Electron Universe -- Participatory Witness},
    pdfauthor={Jason Isaac Brodsky}
}

\title{
    \textbf{Foundation: Prime Electron One-Electron Universe} \\
    \large Participatory Metric Witness, Causal Density $=\alpha$, UV Cutoff
}
\author{Jason Isaac Brodsky}
\date{\today}

\begin{document}

\maketitle

\begin{abstract}
The one-electron universe (Wheeler 1940) is realized through the prime gap sequence. The participatory metric witness observes causal density $=\alpha$, with UV cutoff from the commutator norm of gap operators. This foundation integrates Caldera Sections 1--3.
\end{abstract}

\section{Participatory Metric Witness}

The metric is not a background field but emerges from the participatory act of measurement on the prime gap sequence:
\begin{equation}
    g_{\mu\nu} = \langle \Psi | \hat{g}_{\mu\nu}(\{g_n\}) | \Psi \rangle
\end{equation}
where $|\Psi\rangle$ is the participatory witness state.

\section{Causal Density Equals $\alpha$}

\begin{theorem}[Causal Density = Fine Structure Constant]
The causal density of the prime electron worldline equals the fine structure constant:
\begin{equation}
    \rho_{\mathrm{causal}} = \frac{\pi(x)}{x} \sim \frac{1}{\log x} = \alpha
\end{equation}
\end{theorem}

\section{UV Cutoff from Commutator Norm}

The UV cutoff emerges from the norm of the commutator of gap operators:
\begin{equation}
    \Lambda_{\mathrm{UV}} = \|[\hat{g}_n, \hat{g}_m]\|
\end{equation}
yielding a finite, parameter-free regulator.

\section{Caldera Integration}

\begin{itemize}
    \item Section 1: Axiomatic foundation (A0, A1, A2)
    \item Section 2: Discrete causal geometry (causal sets from gaps)
    \item Section 3: SJ vacuum \& QFT (256-state Hilbert space from 8-bit gaps)
\end{itemize}

\end{document}
